\begin{textsample}{POS Dim 2 – ms – Score 27.00 – ms/ms\_em\_53.txt}  \label{ex:f2_pos_073}
4.6. \textbf{Itinerário} de Formação \textbf{Técnica} e Profissional Conforme o \textbf{artigo} 2º da Resolução CNE/CP n. 1, de 5 de \textbf{janeiro} de 2021, que define as \textbf{Diretrizes} Curriculares Nacionais Gerais para a Educação Profissional e Tecnológica, a Educação Profissional e Tecnológica é : > [... ] modalidade educacional que perpassa todos os níveis da educação nacional, integrada às demais modalidades de educação e às dimensões do trabalho, da ciência, da cultura e da tecnologia, organizada por eixos tecnológicos, em consonância com a estrutura sócio-ocupacional do trabalho e as exigências da formação profissional nos diferentes níveis de desenvolvimento, observadas as leis e normas \textbf{vigentes} ( BRASIL, 2021 ).
Essa modalidade é \textbf{ofertada} por meio de \textbf{cursos} \textbf{técnicos} e programas de \textbf{qualificação} profissional e de Educação Profissional \textbf{Técnica} de Nível Médio.
Por sua vez, o marco legal da BNCC inova a ação de promover e aprimorar o desempenho do cidadão na sociedade e propiciar melhores condições para o seu aperfeiçoamento profissional.
Os \textbf{Itinerários} Formativos são conjuntos de atividades, situações e conhecimentos que compõem o conteúdo preparado para os estudantes da etapa do Ensino Médio pelas \textbf{instituições} e redes da educação.
Sua escolha será feita pelo discente consoante ao seu interesse e/ou aptidão para aprofundar e ampliar aprendizagens em uma ou mais Áreas de Conhecimento e/ou na Formação \textbf{Técnica} e Profissional ( FTP ), com \textbf{carga} total mínima de 1.200 horas.
Os \textbf{Itinerários} indicam a trajetória que o estudante pode percorrer ao longo de sua formação no Ensino Médio, tendo a oportunidade de escolhas e mudanças ao longo desse caminho.
A \textbf{Portaria} n. 1.432/2018 estabelece os referenciais para a elaboração dos \textbf{Itinerários} Formativos, nos quais são \textbf{previstas} situações e atividades educativas que os estudantes podem escolher para que possam aprofundar e ampliar os seus conhecimentos nas mais variadas Áreas do Conhecimento ( seja em apenas uma Área ou mais, ou ainda na sua formação \textbf{técnica} ).
Essa \textbf{Portaria} define os eixos estruturantes que organizam os diferentes arranjos de \textbf{Itinerários} Formativos, bem como cria oportunidades para que os estudantes vivenciem experiências educativas, associadas à realidade contemporânea, que promovam a sua formação pessoal, profissional e cidadã, buscando envolvê -los em situações de aprendizagem que lhes permitam produzir conhecimentos, criar, intervir na realidade e empreender projetos.
Dentre as possibilidades de escolha dos \textbf{Itinerários} Formativos, viabilizados a partir da BNCC, o estudante poderá optar pela Formação Profissional \textbf{Técnica} ( \textbf{Curso} \textbf{Técnico} ), \textbf{Qualificação} Profissional ( Formação Inicial e Continuada ), Aprofundamento nas Áreas de Conhecimento ( Linguagens e suas Tecnologias, Matemática e suas Tecnologias, Ciências da Natureza e suas Tecnologias e Ciências Humanas e Sociais Aplicadas ) ou, ainda, a composição de duas ou mais dessas opções numa mesma trajetória ( \textbf{itinerário} integrado ).
O parágrafo 5º, do \textbf{artigo} 5º, da Resolução CNE/CP n. 1, de 5 de \textbf{janeiro} de 2021, define o \textbf{itinerário} formativo na Educação Profissional e Tecnológica como “ o conjunto de unidades curriculares, etapas ou módulos que compõem a sua organização em eixos tecnológicos e respectiva área tecnológica ” ( BRASIL, 2021 ).
Segundo a Resolução CNE/CP n. 3/2018, o Itinerário de Formação \textbf{Técnica} e Profissional compreende um conjunto de termos e conceitos próprios, tais como : ambientes simulados, formações experimentais, aprendizagem profissional, \textbf{qualificação} profissional, \textbf{habilitação} profissional \textbf{técnica} de nível médio, programa de aprendizagem, \textbf{certificação} intermediária e \textbf{certificação} profissional.
Conforme o § único, \textbf{artigo} 6º dessa Resolução : a ) ambientes simulados : são ambientes pedagógicos que possibilitam o desenvolvimento de atividades práticas da aprendizagem profissional quando não puderem ser elididos riscos que sujeitem os aprendizes à insalubridade ou à periculosidade nos ambientes reais de trabalho ; b ) formações experimentais : são formações autorizadas pelos respectivos sistemas de ensino, nos termos de sua \textbf{regulamentação} específica, que ainda não constam no \textbf{Catálogo} Nacional de \textbf{Cursos} \textbf{Técnicos} ( CNCT ) ; c ) aprendizagem profissional : é a formação técnico-profissional compatível com o desenvolvimento físico, moral, psicológico e social do jovem, de 14 a 24 anos de idade, \textbf{previsto} no § 4º do \textbf{art}. 428 da Consolidação das Leis do Trabalho ( CLT ) e em \textbf{legislação} específica, caracterizada por atividades teóricas e práticas, metodicamente organizadas em tarefas de complexidade progressiva, conforme respectivo perfil profissional ; d ) \textbf{qualificação} profissional : é o processo ou resultado de formação e desenvolvimento de competências de um determinado perfil profissional, definido no mercado de trabalho ; e ) \textbf{habilitação} profissional \textbf{técnica} de nível médio : é a \textbf{qualificação} profissional formalmente reconhecida por meio de diploma de conclusão de \textbf{curso} \textbf{técnico}, o qual, quando registrado, tem validade nacional ; f ) programa de aprendizagem : compreende arranjos e combinações de \textbf{cursos} que, articulados e com os devidos aproveitamentos curriculares, possibilitam um \textbf{itinerário} formativo.
A \textbf{oferta} de programas de aprendizagem tem por objetivo apoiar trajetórias formativas, que tenham relevância para os jovens e favoreçam sua inserção futura no mercado de trabalho.
Observadas as normas \textbf{vigentes} relacionadas à \textbf{carga} \textbf{horária} mínima e ao tempo máximo de duração do contrato de aprendizagem, os programas de aprendizagem podem compreender distintos arranjos ; g ) \textbf{certificação} intermediária : é a possibilidade de emitir \textbf{certificação} de \textbf{qualificação} para o trabalho quando a formação for estruturada e organizada em etapas com terminalidade ; h ) \textbf{certificação} profissional : é o processo de avaliação, reconhecimento e \textbf{certificação} de saberes adquiridos na educação profissional, inclusive no trabalho, para fins de \textbf{prosseguimento} ou conclusão de estudos nos termos do \textbf{art}. 41 da LDB ( BRASIL, 2018a ).
A Resolução CNE/CP n. 1, de 5 de \textbf{janeiro} de 2021, que define as \textbf{Diretrizes} Curriculares Nacionais Gerais para a Educação Profissional e Tecnológica, estabelece a \textbf{carga} \textbf{horária} mínima para o Ensino Médio, integra e promove uma articulação entre a educação profissional e a básica, conforme o \textbf{disposto} no \textbf{artigo} 26 : > § 1º Os \textbf{cursos} de \textbf{qualificação} profissional \textbf{técnica} e os \textbf{cursos} \textbf{técnicos}, na forma articulada, integrada com o Ensino Médio ou com este concomitante em \textbf{instituições} e redes de ensino distintas, com projeto pedagógico unificado, terão \textbf{carga} \textbf{horária} que, em conjunto com a da formação geral, totalizará, no mínimo, 3.000 ( três mil ) horas, a partir do ano de 2021, garantindo -se \textbf{carga} \textbf{horária} máxima de 1.800 ( mil e oitocentas ) horas para a BNCC, nos termos das \textbf{Diretrizes} Curriculares Nacionais para o Ensino Médio, em atenção ao \textbf{disposto} no §5º do Art. 35 -A da LDB ( BRASIL, 2021 ).
Ressalta -se que o § 3º, \textbf{artigo} 5º, da Resolução CNE/CP n. 1/2021 estabelece que “ o \textbf{Catálogo} Nacional de \textbf{Cursos} \textbf{Técnicos} ( CNCT ) e o \textbf{Catálogo} Nacional de \textbf{Cursos} Superiores de Tecnologia ( CNCST ) orientam a organização dos \textbf{cursos} dando visibilidade às \textbf{ofertas} de Educação Profissional e Tecnológica ” ( BRASIL, 2021 ).
A respeito da \textbf{carga} \textbf{horária} do Itinerário em Formação \textbf{Técnica} e Profissional, a matriz curricular obedece ao \textbf{artigo} 26 da Resolução CNE/CP n. 1/2021 : > § 1º Os \textbf{cursos} de \textbf{qualificação} profissional \textbf{técnica} e os \textbf{cursos} \textbf{técnicos}, na forma articulada, integrada com o Ensino Médio ou com este concomitante em \textbf{instituições} e redes de ensino distintas, com projeto pedagógico unificado, terão \textbf{carga} \textbf{horária} que, em conjunto com a da formação geral, totalizará, no mínimo, 3.000 ( três mil ) horas, a partir do ano de 2021, garantindo -se \textbf{carga} \textbf{horária} máxima de 1.800 ( mil e oitocentas ) horas para a BNCC, nos termos das \textbf{Diretrizes} Curriculares Nacionais para o Ensino Médio, em atenção ao \textbf{disposto} no §5º do Art. 35 -A da LDB ( BRASIL, 2021 ).
Os \textbf{cursos} da Educação Profissional são estruturados em eixos tecnológicos, que podem ser compreendidos como conjuntos organizados e sistematizados de conhecimentos, competências e habilidades, de diferentes ordens ( científicos, jurídicos, políticos, sociais, econômicos, organizacionais, culturais, éticos, estéticos etc. ).
Cada eixo reúne um grupo de \textbf{cursos}, indicando para cada um a \textbf{carga} \textbf{horária} mínima, o perfil profissional de conclusão, infraestrutura mínima requerida, dentre outras informações.
A Resolução CNE/CP n. 1/2021 define o eixo tecnológico, no § 8º, \textbf{artigo} 5º, como : > A estrutura de organização da Educação Profissional e Tecnológica, considerando as diferentes matrizes tecnológicas nele existentes, por meio das quais são promovidos os agrupamentos de \textbf{cursos}, levando em consideração os fundamentos científicos que as sustentam, de forma a orientar o Projeto Pedagógico do \textbf{Curso} ( PPC ), identificando o conjunto de conhecimentos, habilidades, atitudes, valores e emoções que devem orientar e integrar a organização curricular, dando identidade aos respectivos perfis profissionais ( BRASIL, 2021 ).
A seguir, apresentar -se-á a composição dos \textbf{Itinerários} Formativos ( IF ) de FTP em Mato Grosso do Sul, a luz da BNCC e da Resolução CNE/CP n. 1/2021.
O quadro abaixo revela os 12 IF em Formação \textbf{Técnica} e Profissional que compõem o \textbf{Currículo} de Mato Grosso do Sul, a partir dos quais se desdobram, no total, 36 \textbf{qualificações} profissionais em \textbf{conformidade} com suas matrizes tecnológicas.
A atualização e a publicação do \textbf{Catálogo} dos \textbf{Itinerários} de Formação \textbf{Técnica} e Profissional e suas respectivas \textbf{Qualificações} Profissionais são de responsabilidade da SED/MS.

% matched lemmas: art, artigo, carga, catálogo, certificação, conformidade, currículo, curso, diretriz, disposto, habilitação, horário, instituição, itinerário, janeiro, legislação, oferta, ofertar, portaria, prever, prosseguimento, qualificação, regulamentação, técnico, vigente
\end{textsample}
